\chapter{Related work}
\label{cap:related-work}

\section{Audio representations for music}
\label{sec:representations}

\todo{From hand-crafted descriptors to learned ones, with the siamese focus of
this work stated up front: MFCCs and their pooling, the broader family of
classical acoustic descriptors (Tzanetakis and Cook 2002, \todo{cite:
Tzanetakis2002}: chroma, spectral contrast, centroid, bandwidth, rolloff, ZCR,
RMS, tempo), the single-Gaussian timbre model compared with symmetric KL
divergence (Mandel and Ellis 2005, \todo{cite: Mandel2005}), convolutional
tagging networks, metric learning with contrastive \cite{Hadsell2006} and
triplet \cite{Schroff2015} losses over three kinds of siamese input
(spectrogram, frozen audio embedding, acoustic descriptors), and self-supervised
models such as MERT \cite{Li2023}.}

\section{Classical classification and similarity measures}
\label{sec:classical-techniques}

\todo{Classical classifiers (kNN, logistic regression, SVM, random forest) used
both to classify genre directly and, through their posterior probabilities, as
a representation for graph construction. Classical similarity measures beyond
cosine: Euclidean, Manhattan, Pearson correlation, Mahalanobis with
regularised (Ledoit-Wolf) covariance, and symmetric KL divergence between
Gaussians. This is the ``alternative classification techniques'' half of the
title.}

\section{Evaluating representations: probing benchmarks}
\label{sec:probing-benchmarks}

\todo{Castellon et al. \cite{Castellon2021} and MARBLE \cite{Yuan2023}. Be
precise about what MARBLE is and is not: it is a probing benchmark over frozen
representations, it contains no similarity or retrieval task at all, and it
ships no MFCC baseline; the MFCC reference numbers come from
\cite{Castellon2021}. See \texttt{docs/ground-truth-study.md}.}

\section{Music similarity and its ground truth}
\label{sec:similarity-ground-truth}

\todo{The odd-one-out triplet game of TagATune \cite{Law2009}, the constraint
derivation and the comparison of adaptation methods of Wolff et al.
\cite{Wolff2012}, the audio-word results of Karamanolakis et al.
\cite{Karamanolakis2016}, the multidimensional judgements of Lee et al.
\cite{Lee2020}, and artist-based relevance on FMA \cite{Cleveland2020}.}

\section{Similarity graphs and hubness}
\label{sec:graphs-hubness}

\todo{Nearest-neighbour graphs, mutual kNN, the hubness phenomenon in
high-dimensional spaces and mutual proximity as a remedy, community detection
and what modularity does and does not say about a music catalogue. Comparing
representations directly: linear CKA (Kornblith et al. 2019, \todo{cite:
Kornblith2019}) as a way to relate representation similarity to graph
similarity, used in the noise-floor and dose-response analysis of SQ3.}

\section{Fuzzy modelling of musical relationships}
\label{sec:fuzzy-related-work}

\todo{Fuzzy set theory (Zadeh 1965, \todo{cite: Zadeh1965}) and fuzzy graphs
(Rosenfeld 1975, \todo{cite: Rosenfeld1975}) as a way to model graded
similarity instead of a crisp yes/no neighbourhood; fuzzy simplicial sets as
used by UMAP (McInnes et al. 2018, \todo{cite: McInnes2018}) as a related
construction. Fuzzy inference systems, Mamdani (Mamdani and Assilian 1975,
\todo{cite: Mamdani1975}) and Takagi-Sugeno (Takagi and Sugeno 1985, \todo{cite:
Takagi1985}), for interpretable, rule-based similarity. Fuzzy k-nearest-neighbour
classification (Keller, Gray and Givens 1985, \todo{cite: Keller1985}) and fuzzy
c-means clustering (Bezdek 1981, \todo{cite: Bezdek1981}) for graded genre
membership, evaluated against multi-label ground truth with the fuzzy
Rand index (Hüllermeier et al. 2012, \todo{cite: Hullermeier2012}). State that
every fuzzy reference is read in full and promoted from \texttt{[to read]} in
\texttt{docs/sources.md} before the report leans on it (V.11).}

\section{Positioning}
\label{sec:positioning}

\todo{One paragraph saying what this work adds: it puts siamese networks with
three kinds of input, classical classification and similarity techniques, and
a fuzzy block, side by side on the same embeddings and the same graphs, at both
the comparative and the absolute level, and it measures the genre proxy against
human judgements instead of assuming it.}
